% =====================================================================
% Supplement（补充材料，独立编译版）：用于期刊主线之外的
% 附加证明、附加模拟、附加数据集分析
% 若目标期刊不区分 appendix/supplement，可删除本文件并在 main.tex 中合并
% =====================================================================
\documentclass[11pt]{article}
\usepackage{amsmath,amssymb,amsthm}
\usepackage{booktabs,graphicx}
\usepackage[margin=1in]{geometry}
\usepackage[colorlinks=true,linkcolor=blue,citecolor=blue]{hyperref}
\usepackage{natbib}

% 与 main.tex 保持一致的定理计数（补充材料中定理编号前缀 S）
\newtheorem{sthm}{Theorem}
\renewcommand{\thesthm}{S\arabic{sthm}}
\newtheorem{slem}[sthm]{Lemma}
\renewcommand{\theslem}{S\arabic{sthm}}

\title{Supplement to ``TODO: 主文标题''}
\author{作者}
\date{}

\begin{document}
\maketitle

\tableofcontents

\section{Additional theoretical results}
% S1 定理及其证明；主文定理依赖的引理若篇幅过长，亦可放此。

\section{Additional simulation studies}
% 补充场景表：建议主文放代表性 2-3 场景，其余全量放补充材料。
% 每张表需在复现包中有对应脚本（sim/文件夹，随机种子固定）。

\section{Additional real data analysis}
% 第二个数据集的全部分析（统计论文顶刊常要求至少两个真实数据集或一个数据集+领域佐证）。

\section{Reproducibility statement}
% 复现包结构声明（TIER 4.0 风格）：
% data/      原始数据（或获取说明与哈希校验）
% code/      全部脚本，随机种子固定，requirements.txt/environment.yml 锁定环境
% output/    一键重新生成的全部图表
% 运行方式、预计时长、软硬件配置逐条列出。

\bibliographystyle{chicago}
\bibliography{refs}
\end{document}
